\chapter{Dei mislukka fundamenta}

\begin{motto}
Eit fag med eit fundament grunnlegg seg ikkje på nytt kvart tjuande år.\\ Designfaget har gjort det i hundre.
\end{motto}

Det finst eit motargument mot heile denne boka, og det er så enkelt at det er lett å gå glipp av. Korleis kan eg seie at design manglar eit fundament, når faget er stappfullt av fundament? Bauhaus bygde ein grunnlære. Ulm bygde ein metode. Simon bygde ein vitskap om det kunstige. Krippendorff gav ut tre hundre sider med undertittelen «eit nytt fundament for design». Det finst etiske manifest frå 1964 og kodeksar frå i går. Faget manglar då ikkje desse tinga. Det eig dei i overflod.

Svaret er å snu motargumentet rundt. Faget har ikkje for få fundament. Det har for mange, og overfloden er ikkje rikdom, ho er sjølve sviktsignaturen. Eit fag som hadde funne grunnen sin, ville slutta å leite. At designfaget i staden har annonsert ein ny grunn, ein ny metode, ei ny etikk om og om att i hundre år, er ei avlesing av kor mange gonger grunnen aldri vart lagd. Talet på fundament er talet på hòl. Det er den tanken kapitlet skal vege, og det finst eitt einaste faktum som veg han: same setning, trykt to gonger om det same fråvêret, tjue år imellom. Me kjem dit til slutt.

\section{Same rørsla, tre stader}

Følg rekkja, men ikkje som ei liste over namn. Følg ho som éi rørsle sett att og att.

Bauhaus opna i Weimar i 1919 med ambisjonen om ein systematisk grunnlære om form. \emph{Vorkurs}, grunnkurset, skulle gje studenten eit objektivt utgangspunkt før verkstaden, ei lære om farge, flate og materiale som låg under all god formgjeving \autocite{gropius1919, itten1963}. Det vart ein pedagogikk og ein stil, og det reiste verda rundt utan grunnlaget med seg. Tretti år seinare opna Ulm med same ambisjonen, no vend mot vitskapen: semiotikk, ergonomi, metode skulle gje faget det Bauhaus ikkje fekk \autocite{maldonado1958, spitz2002ulm}. Skulen kollapsa i 1968, og metodeoptimismen med han.

Resten passerer fortare, fordi rørsla alt er kjend. Metoderørsla bygde ein eksplisitt, rasjonell metode, og hovudmennene forkasta han sjølve \autocite{jones1970methods, alexander1964notes}. Simon gav design ein vitskap om det kunstige; informatikken og økonomien tok imot han, designskulane ikkje \autocite{simon1969sciences}. Cross omdanna draumen frå vitskap til \emph{disiplin} \autocite{cross2001discipline, cross2006}. Nelson og Stolterman la fram «designvegen», ein heil tredje kunnskapskultur \autocite{nelson2012designway}. Blessing og Chakrabarti gav faget ein forskingsmetodologi \autocite{blessing2009drm}. Redström skreiv om korleis design lagar si eiga teori \autocite{redstrom2017making}. Ti fundament, kvart annonsert med dei same orda, \emph{endeleg ein grunn}, og kvart etterfølgt av det neste, som ikkje ville trongst om det førre hadde halde.

Etikken gjer same figuren. I 1964 skreiv Ken Garland og tjueein andre under på First Things First, eit manifest som kravde at designarar vende evnene mot det som tener samfunnet. Det vart fornya i 2000, av ein ny generasjon, fordi det fyrste ikkje hadde endra noko \autocite{firstthings2000}. Profesjonsorgana skreiv kodeksar; det internasjonale designrådet gav ut ein modellkodeks \autocite{icod_code}. Design justice-rørsla la fram ti prinsipp, igjen som om etikken måtte byrje frå null \autocite{costanzachock2020}. Kvar generasjon må \emph{på nytt} minne faget om at det har eit ansvar, fordi minninga frå sist aldri vart bygd inn.

Og språket. Produktsemantikken skulle gje orda for kva form tyder \autocite{krippendorff1984semantics}. Mønsterspråket skulle gje namngjevne, delte løysingar \autocite{alexander1977pattern}. «Intermediate-level knowledge» skulle namngje registeret mellom teori og artefakt \autocite{lowgren2013intermediate}. Archer skulle gjere design til ein disiplin med eige språk \autocite{archer1979discipline}. Kvart forsøk gav nokre nyttige ord; ingen gav faget eit delt vokabular der orda tyder det same for alle, slik «infarkt» tyder det same for to legar som aldri har møttest.

Fundament, etikk, språk er ikkje tre rekkjer. Det er éi rørsle sett tre stader: ein annonsering, eit fråfall, ei ny annonsering. Spørsmålet er kvifor rørsla aldri stansar, og svaret er heile tesen i denne boka i komprimert form. Eit fundament held når det ligg noko under det som tvingar alle til same svar, og som difor ber seg sjølv utan å måtte annonserast på nytt. Designfaget sine fundament er alle \emph{formuleringar}, kvasse og kloke og innflytelsesrike, utan noko prøvbart under seg som kunne gjere usemje om til framgang. Difor forblir dei meiningar. Og ei meining må gjentakast for å overleve, medan kunnskap berre treng å byggjast vidare på.

Dette gjer ikkje design til ein halvferdig kjemi. Kjemien har ein gjenstand som veg likt for kven som helst; biologien har ei art som arvar seg sjølv. Design har ingen slik gjenstand, og det er ikkje ein mangel design kan rette ved å skjerpe seg. Det er kva slag praksis dette er.

\section{Dei formelle som var formelle nok}

Her melder ein motinnvending seg, og ho er sterk. Greitt, dei mjuke fundamenta var formuleringar utan noko prøvbart under. Men faget fostra òg \emph{harde} framlegg, med matematikk og notasjon. Tok ikkje dei heller? Nei. Og at dei ikkje tok, er meir avslørande enn alt det mjuke, for her kan ingen skulde på mangel på stringens. Stringensen var til stades. Framlegget rann ut likevel.

Ta formgrammatikken fyrst. Han vart ikkje født på ein designkonferanse. I 1971 la George Stiny og James Gips fram papiret «Shape grammars and the generative specification of painting and sculpture» på IFIP-kongressen, ein verdskongress for databehandling, og det stod på trykk i \textit{Information Processing 71} året etter \autocite{stiny1972}. Eit formelt system, lånt frå den generative grammatikken i lingvistikken: former vert frambrakte ved reglar som bytar ut ein delform med ein annan, steg for steg, slik ei setning vert frambrakt av omskrivingsreglar. Dette var ekte formalisme, og han let seg implementere på maskin. Stiny bygde han ut over tiår, heilt til boka \textit{Shape} i 2006, der han hevda at det å sjå og rekne med former kunne setjast på heilt eksakt fot \autocite{stiny2006shape}. Apparatet greidde det reint tekniske kunststykket å generere alle villaplanane i ein palladiansk grammatikk. Femti år etter er formgrammatikken ein spesialitet med eigne konferansar og eit dusin trufaste. Han vart fødd mellom informatikarar og enda som ein informatikk-nisje, og aldri den grunnen design reiste seg på. For grammatikken kan generere tusen palladianske villaer; han kan ikkje seie kva for ei av dei som er god. Den dommen sit framleis i auget.

Det same hende med det mest ambisiøse formelle framlegget i nyare tid. Kring tusenårsskiftet utvikla Armand Hatchuel og Benoît Weil ved École des Mines i Paris C-K-teorien: ei matematisk formulert framstilling av design som ei veksling mellom eit konseptrom (C) og eit kunnskapsrom (K), der nyskaping skjer ved at noko ubestemt i C vert gjort bestemt mot K \autocite{hatchuel2009}. Mengdeteoretisk apparat, teorem, undervist i ingeniørmiljø. Men sjå kva all formalismen faktisk seier. C-K skildrar \emph{strukturen} i ein prosess, gangen frå det udefinerte mot det definerte, på eit abstraksjonsnivå så høgt at det er sant for nær sagt all problemløysing, og difor stille om kva som skil ei god form frå ei dårleg der det gjeld. Formalismen er reell, men tom på rett stad. Stiny gjev deg ein generator som spyttar ut gyldige former; Hatchuel gjev deg eit prosesskart over korleis det ukjende vert til det kjende. Begge stoggar i det same augneblinket: når eit menneske må seie at \emph{denne} forma er betre enn dei andre maskinen tillèt.

Den motinnvendinga som freistar her, er at faget berre trong å vere endå hardare, endå meir matematisk, så ville grunnen vore lagd. Stiny og Hatchuel motbeviser ho. Dei var matematiske, dei var harde, og framlegga deira tok like lite som Krippendorffs meining. Men slutninga er ikkje at design feila som ingeniørdisiplin. Formdomen er ikkje eit \emph{knowing that} som ventar på rett notasjon; han er eit \emph{knowing how}, ein taus kunnskap som ligg i ei øvd hand og eit øvd auge, og som difor ikkje let seg redusere til ein syntaks. Å krevje at han skal det, er å krevje at design skal vere ein vitskap det ikkje er. Dei formelle framlegga høyrer difor heime i den same rekkja som dei mjuke, som rekkja sitt sterkaste bevis: sjølv då faget gjorde alt rett etter eiga oppskrift og vart så formelt som råd, fall framlegget i den same grøfta. Det er ikkje måten fundamenta vart lagde på som svikta. Det fanst ingen botn å leggje dei på.

\section{Same setning, trykt to gonger}

Det finst eitt augneblink der påstanden om eit nytt fundament ikkje ligg mellom linene, men står trykt på omslaget. I 2006 gav CRC Press ut ein tjukk, raudbrun bind på vel tre hundre sider. Tittelen var \textit{The Semantic Turn}. Men påstanden låg ikkje i tittelen. Han stod med mindre typar under: «A New Foundation for Design». Eit \emph{nytt} fundament. Orda føreset, utan å seie det, at det gamle ikkje heldt.

Mannen som skreiv det, var Klaus Krippendorff, sytti år gamal og med eit halvt hundreår i akademia bak seg. Han var ikkje kven som helst. Han hadde gått på Ulm, sjølve høgborga for designvitskapen, og seinare, ved Annenberg-skulen ved University of Pennsylvania, bygd ein metode for innhaldsanalyse så stram at reliabilitetskoeffisienten ber namnet hans: Krippendorffs alfa, eit tal to framande forskarar kan rekne ut kvar for seg og få likt. Han visste betre enn nokon i feltet kva ein etterprøvbar grunn var. Og likevel sette denne mannen seg ned, på terskelen til alderdomen, og skreiv tre hundre sider for å reise eit fundament av \emph{meining}. Mannen med sin eigen etterprøvbare koeffisient la ikkje ned ein einaste etterprøvbar stein. Det er kapitlets reinaste ironi, og han er ikkje retorisk: han følgjer av at meining ikkje kan vegast.

Men det avslørande er ikkje boka. Det er at ho ikkje var åleine om setninga. Tjue år tidlegare, i 1986, gav Terry Winograd og Fernando Flores ut \textit{Understanding Computers and Cognition} med same undertittel, ord for ord: «A New Foundation for Design» \autocite{winograd1986}. Dei to bøkene deler ingenting elles. Winograd og Flores byggjer på Heidegger og Maturana og hevdar at god programvare må rekne med korleis menneske er kasta inn i verda; Krippendorff byggjer på semantikk og hevdar at design er meiningsskaping. Ulike fag, tjue år imellom, ingen kontakt. Det einaste dei deler, er løftet på omslaget og fråvêret løftet svarar på. Eit fag der to forskarar uavhengig kan trykkje «eit nytt fundament for design», med dei eksakt same orda, utan at den fyrste boka gjorde den andre overflødig, er eit fag der fundamentet er eit varig hòl i marknaden. Same tittel kan seljast på nytt, fordi det førre salet aldri leverte varen.

Set dette mot eit fag med ein annan slags gjenstand. Kjemien la grunnen sin éin gong, og me kan setje datoen nær på året. Kring 1789 vog Antoine Lavoisier stoff før og etter forbrenning på ei vekt så nøyaktig at ho fanga tap og vinst på milligrammet, og viste at massen er den same på begge sider. To år før hadde han rydda i namna og gjeve grunnstoffa ein nomenklatur. Lavoisier sjølv vart halshoggen 8. mai 1794, men grunnen hans døydde ikkje med han, for han låg ikkje i munnen hans. Han låg i vekta. Eit anna menneske kunne gjenta vegingane og få same svaret, og difor trong ingen seinare skrive «eit nytt fundament for kjemien».

Det er den prøva designfaget aldri står. Krippendorffs meining kunne ikkje vegast. To framande som las same gjenstanden, kunne ikkje balansere ei likning og lande på same svaret; tydinga låg hjå kvar av dei. Difor måtte boka, som alle bøkene før, halde fast på opphavsmannen for å bestå. Etter Lavoisier kom det ingen ny grunn, fordi den fyrste heldt. Etter Winograd kom det ein ny, og etter Krippendorff endå fleire, fordi ingen heldt. Kjemien hadde noko som tvinga alle til same svar, og slutta å leite. Design har inga veging, og difor ein uendeleg serie undertitlar.

\vignett

Den siste i rekkja er den boka her peikar mot, og eg skal ikkje løyne at han er nett det: endå eit forsøk. Men fyrst må klagemålet førast heilt fram, gjennom skadane, fråvêret og avleggjaren, så lesaren kan vege forsøket mot alt det skal bøte på.

\laeringsmal{\begin{itemize}
\item Forklare argumentet om at overfloden av fundament er sjølve avlesinga av fråvêret, og kvifor eit fag berre grunnlegg seg serielt dersom ingen av grunnleggingane heldt.
\item Skilje mellom mjuke fundament (formuleringar) og harde, formelle framlegg, og vise at båe fall i den same grøfta.
\item Bruke formgrammatikken og C-K-teorien til å vise at stringens utan eit objekt å vere stringent om, ikkje gjev eit fundament, og at formdomen er eit \emph{knowing how} som ikkje let seg formalisere.
\item Lese den doble undertittelen hjå Winograd (1986) og Krippendorff (2006) som eit etterprøvbart faktum om at grunnen aldri vart lagd.
\end{itemize}}

\ovingar{\begin{enumerate}
\item Sett opp ei tidsline med dei ti fundamenta i kapitlet, med årstal og dei orda kvart vart annonsert med. Kva er det retoriske fellestrekket, og kva avslører gjentakinga?
\item Vel eit fag som faktisk grunnla seg éin gong (kjemi, biologi). Kva var det under som tvinga alle til same svar? Design manglar ein slik parallell, ikkje fordi det kom for kort, men fordi gjenstanden er av eit anna slag: kva slag?
\item Studer eit konkret framlegg om formgrammatikk eller C-K-teori. Finn nøyaktig punktet der formalismen overlèt avgjerda til ein menneskeleg dom. Kunne det punktet fyllast utan å gjere design til noko anna enn design?
\item Diskuter i seminar: er påstanden «faget har for mange fundament, ikkje for få» rettferdig, eller forvekslar han eit ungt fags utforsking med eit gammalt fags svikt? Argumenter begge vegar.
\end{enumerate}}

\vidarelesing{\fullcite{krippendorff2006}; \fullcite{nelson2012designway}; \fullcite{cross2001discipline}; \fullcite{simon1969sciences}; \fullcite{firstthings2000}; \fullcite{redstrom2017making}.}
